\documentclass[aspectratio=169,11pt]{beamer}

% --- Theme ---
\usetheme{metropolis}
\usepackage{appendixnumberbeamer}

% --- Packages ---
\usepackage{fontspec}
\usepackage{minted}
\setminted{fontsize=\small,breaklines=true}
\usepackage{booktabs}
\usepackage{tikz}
\usepackage{fontawesome5}

% --- Colors ---
\definecolor{healthgreen}{RGB}{34,139,34}
\definecolor{alertred}{RGB}{220,53,69}
\setbeamercolor{alerted text}{fg=alertred}

% --- Metadata ---
\title{Module 4: Descriptive statistics and epidemiological measures}
\subtitle{Data analysis with Python for health specialists}
\author{Ya\'{e} Ulrich Gaba}
\institute{AIRINA Labs}
\date{2026}

\begin{document}

% ============================================================
\maketitle

% ============================================================
\begin{frame}[fragile]{Mean vs.\ median: why it matters}
\begin{minted}{python}
import pandas as pd
glucose = pd.Series([92, 98, 104, 95, 88, 110, 97, 340, 101, 93])

print(f"Mean:   {glucose.mean():.1f} mg/dL")   # 121.8
print(f"Median: {glucose.median():.1f} mg/dL")  # 97.5
\end{minted}

\vspace{6pt}
One patient with DKA (glucose = 340) pulls the mean up by 25~mg/dL. The median is unaffected.

\vspace{6pt}
\begin{alertblock}{Rule of thumb}
Use \textbf{median} for skewed data (lab values, costs, length of stay). Use \textbf{mean} for symmetric data (height, diastolic BP in healthy adults).
\end{alertblock}
\end{frame}

% ============================================================
\begin{frame}[fragile]{Spread: SD, IQR, and coefficient of variation}
\begin{minted}{python}
# Two clinics with similar means but different variability
clinic_a = pd.Series([120, 122, 118, 125, 121, 119, 123, 120])
clinic_b = pd.Series([98, 145, 110, 155, 102, 160, 115, 135])

print(f"Clinic A: mean={clinic_a.mean():.1f}, "
      f"SD={clinic_a.std():.1f}")
print(f"Clinic B: mean={clinic_b.mean():.1f}, "
      f"SD={clinic_b.std():.1f}")
\end{minted}

\vspace{4pt}
Both means $\approx$121, but Clinic~B has patients with \alert{dangerously high and unusually low} blood pressure.

\vspace{4pt}
\textbf{IQR} = Q75 -- Q25. Robust to outliers, like the median.\\
\textbf{CV} = SD/mean $\times$ 100\%. Compares variability across different scales.
\end{frame}

% ============================================================
\begin{frame}[fragile]{Frequency tables and cross-tabulations}
\begin{minted}{python}
# One-way frequency table
freq = recent["le_cat"].value_counts().sort_index()

# Two-way cross-tabulation (the "Table 1" of papers)
ct = pd.crosstab(recent["continent"], recent["le_cat"],
                 margins=True)

# With row percentages
ct_pct = pd.crosstab(recent["continent"], recent["le_cat"],
                     normalize="index").round(3) * 100
\end{minted}

\vspace{4pt}
\begin{exampleblock}{Clinical context}
Every ``Table~1'' in a clinical paper is essentially a cross-tabulation. In pandas, \texttt{pd.crosstab()} builds them in one line.
\end{exampleblock}
\end{frame}

% ============================================================
\begin{frame}[fragile]{Prevalence and incidence}
\begin{columns}
\column{0.48\textwidth}
\textbf{Prevalence} (burden):
$$\frac{\text{Existing cases}}{\text{Total population}} \times 100$$

\begin{minted}{python}
prev = 625 / 5000 * 100
# 12.5%
\end{minted}

\column{0.48\textwidth}
\textbf{Incidence rate} (risk):
$$\frac{\text{New cases}}{\text{Person-time at risk}}$$

\begin{minted}{python}
ir = 18 / 850  # per person-year
# = 21.2 per 1,000 PY
\end{minted}
\end{columns}

\vspace{8pt}
\begin{alertblock}{Key difference}
Low incidence + high prevalence = patients survive a long time (diabetes). High incidence + low prevalence = disease kills quickly (Ebola).
\end{alertblock}
\end{frame}

% ============================================================
\begin{frame}[fragile]{Risk ratio (relative risk)}
Compares probability of disease between exposed and unexposed:
$$\text{RR} = \frac{a/(a+b)}{c/(c+d)}$$

\begin{minted}{python}
# Smoking and lung cancer
a, b, c, d = 80, 920, 10, 990
risk_exposed = a / (a + b)     # 8.0%
risk_unexposed = c / (c + d)   # 1.0%
rr = risk_exposed / risk_unexposed  # 7.92

# 95% CI via log(RR)
import math
se = math.sqrt(1/a - 1/(a+b) + 1/c - 1/(c+d))
ci = (math.exp(math.log(rr) - 1.96*se),
      math.exp(math.log(rr) + 1.96*se))
\end{minted}

Smokers have \alert{8$\times$ the risk} of lung cancer compared to non-smokers.
\end{frame}

% ============================================================
\begin{frame}[fragile]{Odds ratio}
Used in \textbf{case-control} studies where you cannot compute risk directly:
$$\text{OR} = \frac{a \times d}{b \times c}$$

\begin{minted}{python}
# Obesity and knee osteoarthritis
a, b, c, d = 120, 60, 80, 140
odds_ratio = (a * d) / (b * c)  # 3.50

se = math.sqrt(1/a + 1/b + 1/c + 1/d)
ci = (math.exp(math.log(odds_ratio) - 1.96*se),
      math.exp(math.log(odds_ratio) + 1.96*se))
\end{minted}

\vspace{4pt}
\begin{exampleblock}{When OR $\approx$ RR}
The OR approximates the RR when the disease is \textbf{rare} ($<$10\% prevalence). For common outcomes, OR exaggerates the association.
\end{exampleblock}
\end{frame}

% ============================================================
\begin{frame}[fragile]{Age-standardized rates}
Crude rates mislead when populations have different age structures.

\begin{minted}{python}
# Direct standardization
std_weight = [0.26, 0.42, 0.22, 0.10]  # WHO World Std

# Age-specific rates x standard weights
asr_a = (data_a["rate"] * std_weight).sum()
asr_b = (data_b["rate"] * std_weight).sum()
\end{minted}

\vspace{4pt}
\begin{alertblock}{Always report your standard}
Different standards (WHO, European, US 2000) give different results. You must use the \textbf{same} standard to compare with published rates.
\end{alertblock}
\end{frame}

% ============================================================
\begin{frame}{What you can do after this module}
\begin{enumerate}
  \item Choose the right measure of center (mean vs.\ median) for health data
  \item Compute prevalence and incidence with confidence intervals
  \item Build 2$\times$2 tables and calculate risk ratios and odds ratios
  \item Construct cross-tabulations for ``Table 1'' style reporting
  \item Perform age standardization to fairly compare populations
\end{enumerate}

\vspace{12pt}
\centering
\Large \alert{Next: Module 5 --- Visualization for health data}
\end{frame}

% ============================================================
\begin{frame}{Exercises (Hour 3)}
\begin{enumerate}
  \item \textbf{Summary stats:} Compute mean, median, SD, IQR for life expectancy by continent (2007)
  \item \textbf{Odds ratio:} From a 2$\times$2 table (200 cases, 200 controls), compute OR with 95\% CI
  \item \textbf{Age standardization:} Compare crude vs.\ standardized mortality for two populations
  \item \textbf{Mini-project:} Build an epidemiological profile using Gapminder + WHO data
\end{enumerate}
\end{frame}

\end{document}
